\label{ch:results}

This chapter presents, interprets, and summarises the real numerical
results of \texttt{inference5\_icecube.ipynb} (in
\texttt{notebooks/inference/}), executed end to end against the real
IceCube DeepCore 9-year golden-event data (\cref{ch:data-release}) using
the physics of \cref{ch:detector-implementation} and the inference
machinery of \cref{ch:inference-process}. Every number below is a real,
reproducible output of that notebook or of this project's own test suite,
not illustrative.

\section{Detector-level closure check}

\begin{resultbox}
Before any fit, a single forward-model evaluation at NuFIT~6.1 defaults
($\nu_{\rm norm}=\mu_{\rm norm}=1$) takes $2.7$~s and predicts a total of
$513.2$ events, of which the muon background alone contributes $512.17$
-- confirming the neutrino signal's own $\sim$0.2\% share at this
unnatural normalization (\cref{sec:free-normalization-icecube}).
Autograd gradients with respect to all four free parameters are finite:
\begin{equation}
  \frac{\partial(\text{total})}{\partial(\theta_{23},\Delta m^2_{3l},\nu_{\rm norm},\mu_{\rm norm})}
  = (-0.083,\ -65.74,\ 1.049,\ 512.17),
\end{equation}
confirming gradient flow through the real Earth-matter transition
probability, the real Honda-flux-weighted MC reweighting, and the
differentiable per-bin histogram, before any fit is attempted.
\end{resultbox}

\section{Real fit}
\label{sec:results-fit}

\begin{resultbox}
Starting from the calibrated point (\labelcref{eq:calibration}), the real
4-parameter Poisson fit converges in 33 closure evaluations
($131$~s wall time):
\begin{equation}
  \keyresult{\theta_{23} = 45.42^\circ \pm 1.85^\circ\ (\sin^2\theta_{23}=0.5074), \qquad
    \Delta m^2_{3l} = (2.3947\pm0.0392)\times10^{-3}\,{\rm eV}^2,}
\end{equation}
\begin{equation}
  \nu_{\rm norm} = 19{,}977.0\pm0.4, \qquad \mu_{\rm norm} = 1.617\pm0.216,
\end{equation}
with $-2\ln L$ improving from $353.71$ to $340.44$.
\Cref{tab:icecube-comparison} compares the oscillation-parameter result
directly against IceCube's own published best fit: the pull on
$\sin^2\theta_{23}$ is $-0.05\sigma$ -- close agreement, not exact
identity, exactly as expected given this document's own fixed-systematics
scope decision (\cref{sec:flux-hypersurface}).
\end{resultbox}

\begin{table}[htb]
\centering
\caption{This document's own fit vs.\ IceCube's own published best fit (Table III/IV of \cite{Abbasi2023IceCube}).}
\label{tab:icecube-comparison}
\begin{tabular}{@{}lrr@{}}
\toprule
Parameter & \texttt{tpeanuts} (this document) & IceCube published \\
\midrule
$\sin^2\theta_{23}$ & $0.5074\pm0.033$ & $0.51$ \\
$\Delta m^2_{32}$ [$10^{-3}$ eV$^2$] & $2.32\pm0.04$ & $2.41\pm0.07$ \\
Atmospheric-muon norm.\ scale & $1.617\pm0.216$ & $1.39$ \\
\bottomrule
\end{tabular}
\end{table}

\begin{notebox}
Uncertainties on $\sin^2\theta_{23}$ and $\Delta m^2_{32}$ are propagated
from $\theta_{23}$'s/$\Delta m^2_{3l}$'s own Laplace $\sigma$
(\cref{sec:results-fit}) via $d(\sin^2\theta_{23})/d\theta_{23}$;
$\Delta m^2_{32}=\Delta m^2_{3l}-\Delta m^2_{21}$, this project's own
convention (\cref{ch:detector-implementation}).
\end{notebox}

\begin{figure}[htb]
\centering
\includegraphics[width=0.75\textwidth]{fig1_fit_history}
\caption{Poisson $-2\ln L$ at every LBFGS closure evaluation (including internal line-search evaluations).}
\label{fig:fit-history}
\end{figure}

\section{2-D likelihood slice}

\begin{figure}[htb]
\centering
\includegraphics[width=0.62\textwidth]{fig2_2d_scan}
\caption{$-2\ln L$ over $(\theta_{23},\Delta m^2_{3l})$, $\nu_{\rm norm}$/$\mu_{\rm norm}$ fixed at their own best-fit values ($625$ forward evaluations, $1028$~s).}
\label{fig:2d-scan}
\end{figure}

\begin{notebox}
\Cref{fig:2d-scan} is a \textbf{fixed-normalization slice}, not a fully
profiled contour: $\nu_{\rm norm}$/$\mu_{\rm norm}$ are held at
\cref{sec:results-fit}'s own best-fit values at every grid point, one
forward pass each, rather than re-minimized (the \pyfunc{loglik\_grid}
machinery documented in the SNO-II companion document's own Chapter~4
supports \code{profile\_others=True} directly and could be applied here
unchanged -- an explicitly deferred next step, not attempted in this
notebook).
\end{notebox}

\section{Real data comparison}

\begin{figure}[htb]
\centering
\includegraphics[width=0.85\textwidth]{fig3_real_comparison}
\caption{Real IceCube 200-bin observed spectrum vs.\ two \texttt{tpeanuts} curves (NuFIT~6.1 simulation and the fitted result), collapsed onto reconstructed energy.}
\label{fig:real-comparison}
\end{figure}

\section{The muon-background normalization audit}
\label{sec:results-munorm}

\begin{resultbox}
\Cref{sec:munorm-audit}'s audit, carried out quantitatively:
\begin{itemize}
  \item \textbf{Earth-matter unitarity, ruled out.} At the published
    best-fit point, forcing \code{reunitarize=True} changes the total
    predicted rate from $21{,}913.81$ to $21{,}913.37$ (of a real observed
    $21{,}914.0$) -- a $0.002\%$ shift, with a maximum per-bin relative
    difference of only $0.06\%$. Three orders of magnitude too small to
    account for $\mu_{\rm norm}$'s own $16\%$ gap above the published
    $1.39$.
  \item \textbf{Fixed detector systematics, identified and quantified.}
    The real hypersurface correction's own mean value (best-fit vs.\
    nominal systematics point) shifts by $+2.0\%$ (\code{nc}, identical to
    \code{nue\_cc}), $+2.0\%$ (\code{numu\_cc}), and $+12.0\%$
    (\code{nutau\_cc}) -- a real, non-negligible normalization effect that
    this project's own reduced (4-parameter, no-prior) fit cannot
    distinguish from a genuine muon-background rescaling, unlike
    IceCube's own $\sim$40-parameter joint fit with Gaussian priors on
    every nuisance.
\end{itemize}
\end{resultbox}

\section{Summary and known limitations}

This document fits IceCube's real public golden-event data using its own
event-by-event weighted-Monte-Carlo reweighting technique
(\cref{ch:detector-implementation}) and a real Poisson likelihood
(\cref{ch:inference-process}), recovering $\theta_{23}$/$\Delta m^2_{3l}$
close to the collaboration's own published best fit ($-0.05\sigma$ pull on
$\sin^2\theta_{23}$). A dedicated test suite
(\cref{sec:test-suite}), the first for this project's only weighted-MC
detector, and a quantitative audit of the fitted muon-background
normalization's $16\%$ discrepancy against IceCube's own published value
(\cref{sec:munorm-audit}) -- ruling out an Earth-matter unitarity
explanation and identifying fixed detector systematics as the more likely
cause, with the size of the effect quantified directly -- are this
document's own contributions beyond reproducing the existing notebook.

\begin{notebox}
\textbf{Known limitations, stated explicitly.} The 5 real
detector-systematics parameters are fixed at their published best-fit
point rather than re-fit as free nuisances -- the natural next step the
muon-norm audit itself motivates, and the real hypersurface machinery
(\cref{sec:flux-hypersurface}) already supports it without further code
changes. The real MC is downsampled $10\times$ for computational
tractability. Production is treated at the top of the atmosphere rather
than through a full atmosphere-production-height treatment (a
sub-percent-level simplification). \Cref{fig:2d-scan}'s likelihood
picture is a fixed-normalization slice, not a fully profiled contour --
the SNO-II companion document's own \code{profile\_others=True} machinery
could be applied here directly. Next steps: free the 5 systematics as
real Gaussian-prior nuisances and re-audit $\mu_{\rm norm}$; profile
$\nu_{\rm norm}$/$\mu_{\rm norm}$ at every 2-D scan grid point; and
combine this real atmospheric measurement with this project's other real
detectors' oscillation-parameter constraints in a genuine joint
likelihood.
\end{notebox}
